% Generated by src/targets.py -- do not edit by hand.
% Nominal-proxy table (referee comment 10).  Intervals are bounded ranges
% over the stated albedo and radius-factor assumptions, NOT confidence
% levels; rows marked (a) have M sin i above the Chen & Kipping (2017)
% Terran/Neptunian transition and use a rocky extrapolation.
\begin{table*}[t]
\centering
\caption{\rev{Nearest confirmed temperate planets (plus one solar-type RV
signal whose planet interpretation has since been refuted) as SGL imaging
targets -- explicitly a \emph{nominal-proxy} 
table, not a propagated uncertainty budget.  Radii use the Chen \&
Kipping (2017) Terran branch $R = 1.01\,(M\sin i/M_\oplus)^{0.279}
\,R_\oplus$ evaluated at $M\sin i$ (no de-inclination correction is
folded into the nominal value; the conditional isotropic factors are
given in the text); $\pm$ values are the $4.03\%$ intrinsic dispersion of
that branch.  $\mathrm{SNR_C}$ follows the scalar rule of eq.~\ref{eq:snrrule}
and the bracket is the bounded product of $A\in[0.15,0.45]$ with a
$\pm15\%$ radius allowance -- an assumption, not a confidence interval.
$\Delta v_{90}$ is the time-integral of the projected transverse
acceleration over 90\,d for a face-on orbit, with the range over starting
mean anomaly and the corresponding edge-on value in parentheses.
Observing geometry $z=650$\,AU, $d=1$\,m, $\lambda=1\,\mu$m, 1800\,s.}}
\label{tab:targets}
\scriptsize
\setlength{\tabcolsep}{2.0pt}
{\bfseries % entire table body is revised material (editor: bold changes)
\resizebox{\textwidth}{!}{%
\begin{tabular}{lcccccccccccc}
Planet & Host (SpT) & $z_0$ & $M\sin i$ & $R_p$\,$\pm$\,disp.$^{d}$ & $S$ & $P$ &
$D_{\rm img}$ & $\mathrm{SNR_C}$\,[range]$^{b}$ & Focal direction$^{c}$ & $\beta_{\rm ecl}$ &
$\Delta v_{90}$ (face)$^{b}$ \\
& & (pc) & ($M_\oplus$) & ($R_\oplus$) & ($S_\oplus$) & (d) & (km) &
(1800\,s) & (J2000) & (deg) & ($\mathrm{km\,s^{-1}}$ [range], edge) \\
\midrule
Proxima Cen b & Proxima Centauri (M5.5V) & 1.30 & 1.07 & 1.03\,$\pm$\,0.04 & 0.65 & 11.19 & 31.8 & 666\,[283--1148] & 02$^{\rm h}$30$^{\rm m}$, $+62^{\circ}41'$ & +44.8 & 5.81 (3.70)\\
Ross 128 b & Ross 128 (M4V) & 3.38 & 1.35 & 1.10\,$\pm$\,0.04 & 1.38 & 9.87 & 13.1 & 581\,[247--1003] & 23$^{\rm h}$48$^{\rm m}$, $-00^{\circ}48'$ & +0.5 & 2.94 (1.87)\\
GJ 1061 d & GJ 1061 (M5.5V) & 3.67 & 1.64 & 1.16\,$\pm$\,0.05 & 0.69 & 13.03 & 12.7 & 282\,[120--487] & 15$^{\rm h}$36$^{\rm m}$, $+44^{\circ}31'$ & +60.8 & 1.68 (1.07)\\
Teegarden c & Teegarden's Star (M7V) & 3.83 & 1.11 & 1.04\,$\pm$\,0.04 & 0.37 & 11.41 & 10.9 & 130\,[55--224] & 14$^{\rm h}$53$^{\rm m}$, $-16^{\circ}53'$ & -0.3 & 1.72 (1.10)\\
GJ 273 b\textsuperscript{a} & GJ 273 (M3.5V) & 3.80 & 2.89 & 1.36\,$\pm$\,0.06 & 1.06 & 18.65 & 14.4 & 490\,[208--846] & 19$^{\rm h}$27$^{\rm m}$, $-05^{\circ}14'$ & +16.5 & 1.34 (0.85)\\
$\tau$ Cet e\textsuperscript{a}\textsuperscript{e} & $\tau$ Ceti (G8.5V) & 3.60 & 3.93 & 1.48\,$\pm$\,0.06 & 1.71 & 162.90 & 16.5 & 909\,[386--1568] & 13$^{\rm h}$44$^{\rm m}$, $+15^{\circ}56'$ & +24.8 & 0.11 (0.07)\\
\bottomrule
\end{tabular}
}
}

\vspace{1mm}
\parbox{0.96\textwidth}{\footnotesize
\textbf{$^{a}$}\,$M\sin i$ lies above the Chen \& Kipping (2017) Terran/Neptunian transition ($2.04^{+0.66}_{-0.59}\,M_\oplus$), so the rocky-branch radius used here is an extrapolation, not a measurement.
\textbf{$^{b}$}The $\mathrm{SNR_C}$ bracket is a \emph{bounded range} over the stated albedo ($A\in[0.15,0.45]$) and radius allowance ($\pm15\%$), not a confidence level and not propagated from catalogue uncertainties. The $\Delta v_{90}$ entries give the face-on value with the corresponding edge-on value in parentheses; the spread over the starting mean anomaly inside a 90-day window is below $0.02\,\mathrm{km\,s^{-1}}$ for every row except $\tau$~Cet~e ($0.09$--$0.13\,\mathrm{km\,s^{-1}}$).
\textbf{$^{c}$}Focal direction is the antipodal celestial coordinate of the host; $\beta_{\rm ecl}$ is its ecliptic latitude, which sets the departure plane-change cost.
\textbf{$^{d}$}Radii, stellar parameters and orbital elements are from the references in Section~\ref{sec:targets}; none of these planets transits, so $M\sin i$ (not $M$) is the observed quantity.
\textbf{$^{e}$}\,$\tau$~Cet~e is listed for dynamics illustration only. The NASA Exoplanet Archive carries its 162.9-day signal as a \emph{false positive} after \citet{figueira25}, who attribute the $\tau$~Ceti RV signals to stellar activity; it is therefore not a confirmed target and its radius, $S$, SNR and $\Delta v$ inherit the assumed $M\sin i = 3.93\,M_\oplus$ of a refuted ephemeris.
}
\end{table*}
